\chapter{Test d'utilit\'e FEDS et validation OC5}\label{chap:experts}

% TODO: rédiger après collecte du panel.
% Sources : documentation/managerial_playbook.md, documentation/implications_manageriales_et_pratiques.md
% Sections cibles (cf. plan détaillé §7.1 à §7.6) :
%   7.1 Justification du test d'utilité au sens FEDS (Venable et al. 2016)
%   7.2 Protocole du test d'utilité (panel 5-8 experts, démo, questionnaire Likert structuré sur 4 dimensions)
%   7.3 Résultats quantitatifs Likert [PLACEHOLDER : dépend collecte]
%   7.4 Observations post-démo structurées
%   7.5 Audit trail et conformité Article 14 EU AI Act
%   7.6 Validation OC5

\section{Justification du test d'utilit\'e au sens FEDS}\label{sec:experts-justif}

\textit{\`A r\'ediger.}

\section{Protocole du test d'utilit\'e}\label{sec:experts-protocole}

\textit{\`A r\'ediger.}

\section{R\'esultats quantitatifs Likert}\label{sec:experts-likert}

\textit{Placeholder~: \`a compl\'eter apr\`es collecte du panel.}

\section{Observations post-d\'emo structur\'ees}\label{sec:experts-observations}

\textit{Placeholder~: \`a compl\'eter apr\`es collecte du panel.}

\section{Audit trail et conformit\'e Article~14 EU AI Act}\label{sec:experts-audit}

\textit{\`A r\'ediger.}

\section{Validation OC5}\label{sec:experts-validation}

\textit{\`A r\'ediger.}
